% ============================================================
%  P4 paper skeleton v0.1  ·  BioSNN-Plug Arbitration Paper
%  Build:  latexmk -pdf main  (needs hyperref, booktabs, graphicx, amsmath)
%  写作纪律 (inherited from P2): 数字全部来自 research/ 复现记录;
%  \todo{} 是任务清单; 一个数字只有一个家; 不利结果照登.
% ============================================================
\documentclass[11pt]{article}
\usepackage[margin=1in]{geometry}
\usepackage[T1]{fontenc}
\usepackage[utf8]{inputenc}
\usepackage{booktabs, graphicx, hyperref, amsmath, xcolor}
\newcommand{\todo}[1]{\textcolor{red}{[TODO: #1]}}
\title{When Can Arbitration Be Measured?\\Mixed Local Plasticity Rules Under One Protocol}
\author{Jinyao Shen\thanks{Independent researcher. Correspondence: \todo{GitHub URL}.}\\ \small{github.com/lzmd-arch/BioSNN-Plug}}
\date{November 2026}

\begin{document}
\maketitle

% ------------------------------------------------------------
\begin{abstract}
\todo{150-200 words. Outline: (1) problem --- multiple local learning rules acting on
the same synapses produce conflicting updates; single-rule systems never face this;
(2) what we built --- a continuous-mixing arbitration between heterogeneous local
rules (Hebbian / e-prop / R-STDP), meta-learned by an evolution-strategy outer loop,
under a retention constraint that excludes degenerate solutions, in a network with no
backpropagation anywhere; (3) headline numbers --- \todo{synergy-benchmark result};
\todo{conflict reduction reported jointly with retention under the two-sided constraint $\alpha \in [0.2, 0.8]$}; (4) stance --- arbitration is measured,
not assumed; where it fails is reported.}
\end{abstract}

% ------------------------------------------------------------
\section{Introduction}
\subsection{The problem this paper measures}
% 动机三层: (a) 生物学 - Science 2025: 同一神经元不同突触遵循不同可塑性规则
%  (引用进仓后补); (b) 工程学 - 多层系统中 Hebbian/e-prop/R-STDP 必然共享突触;
%  (c) 认识论 - 单方规则各有信息局限, 混合是否 > 单方是可测问题.
\subsection{Why arbitration is not hyperparameter tuning}
% 划界: 静态学习率调度 vs 随任务状态自适应的混合权重; 一句话立场.
\subsection{Contributions}
\begin{itemize}
  \item \textbf{Formalisation}: arbitration as continuous mixing under a retention constraint;
  degenerate solutions ($\alpha \to 0$) characterised and excluded by protocol
  (inherited from \todo{P2 citation, same venue family}).
  \item \textbf{Synergy benchmark}: \todo{one-line task description where each rule alone
  underperforms the mixture} --- the first benchmark in which arbitration has a
  non-trivial optimum to find.
  \item \textbf{Mechanism comparison}: continuous mixing (ours) vs.\ discrete per-synapse
  rule switching (Turcu et al., 2025) under matched conditions.
  \item \textbf{Open assets}: rules, benchmark, pre-registration records, arbitration code.
\end{itemize}

% ------------------------------------------------------------
\section{Related Work}
\subsection{Switching synaptic plasticity rules}   % Turcu 2025 = closest neighbour.
% 必须 delimit 的两点: 离散 vs 连续; 其强规则含 BP vs 我们全程纯局部 (含外层 ES).
% [2026-10-24 裁决] 实证对比**未做**, 按分析性对比 + 我们的设计理由写;
% 实证对比进 future work, 并明确声明不可直接比较（不同条件, 按 P2 预算纪律处理）.
\subsection{Meta-learned plasticity}               % Confavreux 2025 (ES, single rule family,
% failure mode beyond simple parameterisations -> our 1-scalar deliberate choice);
% Miconi 2018 as contrast (global gradient meta-learning).
\subsection{Neuromodulation and gating}            % Zanardi 2024 meta-plasticity as memory;
% Frémaux multi-modulator framing; 划界: 单规则内调制 vs 规则间仲裁.
\subsection{Multi-objective interference at the synapse}  % EWC/任务级冲突 resolution
% vs 规则级冲突; 一句话划界.

% ------------------------------------------------------------
\section{Methods}
\subsection{Arbitration formalism}
% 两/三规则对同一突触的更新 u_A, u_B; 混合更新 u = (1-a) u_A + a u_B;
% 冲突度量 = per-row cosine (继承 P2: per-synapse 退化已证).
% 约束是**双侧**的: alpha in [0.2, 0.8] —— 两规则各保留至少 20% 权重("都不许被关掉");
% retention 的定义 = 被压制一方在仲裁结果中实际保留的权重, 与冲突降幅**联合报告**
% (继承 P2 triviality analysis: alpha->0 退化解的识别与封堵; alpha->1 是同一枚硬币的另一面).
\subsection{Meta-learning the arbitration signal}
% 外层 ES, 参数化刻意最小 (1 scalar -> per-row vector 的扩展实验);
% Confavreux 失败模式的直接检验; 预注册: 参数化扩展的停止规则.
% [2026-10-24 条件预注册] 标量仲裁若在本基准显示无增益 (R1/R2 结局), per-row 扩展
% **不予开展** —— 组合空间 xN, 但前提已被否定. 这条写死, 防止未来"再试一维"的无限救援.
\subsection{The Rules-Collide benchmark}
% [2026-10-24 裁决] 协同性基准正式命名 **the Rules-Collide benchmark** ——
% evocative 的名字给工件, 精确的名字("When Can Arbitration Be Measured?")给论文.
% 核心节. 设计判据: 单方 < 混合才可解 (complementarity 任务);
% \todo{任务规格: Hebbian 供特征结构 + e-prop 供时序标签的复合任务};
% 负向验证: 单方基线必须如预期失败 (预注册期望值).
\subsection{Triviality analyses: four cases}\label{sec:triviality}
% [2026-10-24 裁决] 正文放**陈述与教训**(每案四五行: 现象 -> 机制 -> 纪律),
% 推导细节与参数放附录 C. 完整草稿见 paper/p4/section-3.3-triviality-draft.md.
% 四案: ① alpha->0 退化; ② 逐突触余弦退化; ③ 整数目标 ES 惰性; ④ 代理-构念错位.
% NOTE 编号: 裁决里称本节为 "§3.3", 那是按当时的草稿结构; 本骨架 Methods 已有四个小节,
% 插在 §3.4 之后即本节编号为 3.5, protocol inheritance 顺延为 3.6 —— 待作者确认编号.

\subsection{Protocol inheritance}
% 从 P2 继承的六件套 (预算纪律/config dump/逐位复现/冻结随机流/验证集选择/联合报告)
% 一句话引用, 不重写.

% ------------------------------------------------------------
\section{Experimental Setup}
% RTX 5060 8GB; 与 P2 相同的软件栈表; 与 P2 W2 相同的 256 神经元骨干 (可复用资产).

% ------------------------------------------------------------
\section{Results}
\subsection{Synergy: does the mixture beat both parents?}\label{sec:synergy}

The Rules-Collide benchmark pits three arms against each other on a fixed-permutation
sequential MNIST with 1\% of the training labels: Hebbian alone (H), e-prop alone (E), and a
50/50 mixture of the two (MIX). All three share the same $28 \to 256$ input weight, receive the
same data and the same batches, and are scored by the Fisher ratio of the representation they
produce, evaluated on a held-out probe set. Table~\ref{tab:synergy} reports five seeds per arm.

\begin{table}[t]
\centering
\caption{Representation separability (Fisher ratio) on the Rules-Collide benchmark,
40 epochs, 1\% labels, five seeds. Intervals are the interquartile range.
The paired synergy margin subtracts, per seed, the stronger of the two single-rule arms.}
\label{tab:synergy}
\begin{tabular}{lccc}
\toprule
Arm & Fisher (median) & IQR & Seeds used \\
\midrule
H alone (pure Hebbian)   & 0.472 & [0.471, 0.472] & 5/5 \\
E alone (pure e-prop)    & 0.331 & [0.327, 0.337] & 4/5 (one seed failed the run gate) \\
MIX (fixed $\alpha = 0.5$) & \textbf{0.651} & [0.636, 0.652] & 5/5 \\
\midrule
Synergy margin (MIX $-$ best single) & \textbf{+0.179} & [+0.175, +0.182] & 4 paired seeds \\
\bottomrule
\end{tabular}
\end{table}

The mixture beats both parents. The pre-registered criterion was a paired median margin of at
least $\delta = 0.05$ together with a separation requirement (the mixture's lower quartile above
the best single arm's upper quartile); both hold, with the margin more than three times the
threshold. One e-prop seed is excluded: its weight moved less than the run gate allows
($\|\Delta W\|/\|W_0\| = 0.0474 < 0.05$), so it is void rather than slow, and the paired margin
is computed on four seeds, which is stated wherever the number appears.

A static sweep of the mixing weight locates the optimum this benchmark allows. Over
$\alpha \in \{0.3, 0.4, 0.5, 0.6, 0.7\}$ with five seeds each, the Fisher ratio is 0.364
[0.361, 0.450] at 0.3, 0.577 [0.563, 0.627] at 0.4, 0.651 [0.636, 0.652] at 0.5, 0.643
[0.639, 0.643] at 0.6 and 0.621 [0.604, 0.622] at 0.7. The optimum is interior and sits at
equal weights; the adjacent $\alpha = 0.6$ is statistically indistinguishable (paired median
$-0.008$), while 0.4 and 0.7 are clearly worse. Two consequences follow. First, the 50/50
default the benchmark shipped with happens to be at the optimum. Second, an adaptive scheme
has essentially no headroom left to beat it on this task, and that is the question
Section~
ef{sec:pending} asks.

What this result does \emph{not} establish is that \emph{arbitration} is what produced it. The
mixture here is a fixed 50/50 blend; the arm that would adapt the mixing weight is treated
separately in Section~\ref{sec:pending}. Table~\ref{tab:synergy} is therefore a two-level
ordering --- mixture over singles --- and the third level is deliberately left open.
\subsection{Arbitration without degeneracy}\label{sec:pending}

The adaptive arm was run under a pre-registration written before any of its numbers existed,
and the answer is negative. We report it in the order the constraints were discovered.

\textbf{First, no objective could be priced.} The natural objective --- an evolution strategy
climbing on the count of conflict events --- turned out to have no gradient: at $\alpha = 0.5$
the blend already conflicts with the other rule zero times, every perturbation scores the same,
and $\alpha$ never moves; the arm was bit-identical to a frozen $0.5$ in all five seeds. Two
replacements were then checked against a pre-registered discrimination rule and failed it: an
end-to-end Fisher objective has no discriminating power on probes short enough to afford (its
range across $\alpha$ is $0.026$--$0.039$, not growing with probe length), and a smooth proxy
for ``agreement with both rules'' peaks at $\alpha = 0.9$ while assigning the same value to
$\alpha = 0$ and $\alpha = 1$ --- arms that differ by $0.14$ in Fisher. Adaptive arbitration is
therefore \emph{not measurable by any objective we could afford}; the only known-valid objective,
a full-course probe, is priced at about seven hours per fit and is left as future work.

\textbf{Second, the static optimum is itself the default.} A sweep of the mixing weight over
$\alpha \in \{0.3, \dots, 0.7\}$ (Section~\ref{sec:synergy}) locates the optimum at equal
weights. A frozen $0.5$ is therefore not a straw man but the best static setting this benchmark
admits.

\textbf{Third, slow adaptation moves, and moves the wrong way.} With no outer optimizer and no
probe, we let $\alpha(t)$ drift at routine validation checkpoints (step $0.05$, interval
$[0.2, 0.8]$, every five epochs, validation carved from the training split so that the choice
never touches the evaluation set). The trajectory is monotone and identical in shape across all
five seeds: $0.50 \to 0.55 \to \dots \to 0.80$, parking at the upper bound. Final Fisher is
$0.642$ against $0.632$ for a matched fixed-$0.5$ control trained on the same split --- a paired
median of $+0.003$, below the pre-registered $\delta = 0.05$ --- and $0.651$ for the full-data
frozen $0.5$, i.e.\ \emph{worse} by $0.009$. The drift tracks a \emph{transient} validation
signal rather than the final optimum, and it ends by parking on the constraint boundary --- the
same family as the two degeneracies of Section~\ref{sec:triviality}: a quantity that is
optimised, but not the quantity that matters.

\textbf{The line is sealed.} The pre-registration allowed these two rescue arms and no third.
Both are negative, so the conclusion is not ``untested'' but \emph{measured and negative}: on
this benchmark no adaptive arbitration scheme we could build beats a fixed 50/50 mixture, and the
best static mixture is itself 50/50. What remains open is narrower and better posed: whether an
objective exists that is affordable \emph{and} discriminating, which we have shown is the binding
constraint rather than the arbitration idea itself.
\subsection{Continuous mixing vs.\ discrete switching}

The closest neighbour switches plasticity rules \emph{discretely}, synapse by synapse
\citep{turcu2025switching}; we mix them \emph{continuously} with one scalar. We did not run a
matched empirical comparison, and we do not claim one: the conditions differ (their strongest
rule includes backpropagation, ours is local end to end including the outer loop), so a side-by-side
number would compare two different claims at once. What we can state analytically is where the two
formulations differ, and why we chose ours.

\textbf{State space.} Discrete switching gives each synapse a categorical state; for $n$ rules the
per-synapse state space is $n^{|W|}$, and the arbitration object is a selector. Continuous mixing
gives each synapse the \emph{same} scalar $\alpha$, so the arbitration object is one number for the
entire network. The second is a much smaller hypothesis class --- a feature for a paper whose
subject is measurement, a limitation for a paper whose subject is capacity.

\textbf{Attenuation.} Switching applies one rule's update in full; mixing applies a convex
combination, so each rule's contribution is scaled by its weight. In our benchmark this shows up
directly: the mixture's weight displacement per run ($\|\Delta W\|/\|W_0\| \approx 0.41$) sits
between pure Hebbian's ($\approx 0.71$) and pure e-prop's ($\approx 0.05$), i.e.\ mixing costs
magnitude relative to the dominant rule. Whether switching's fuller updates are worth the larger
state space is an empirical question this paper does not answer.

\textbf{Searchability.} A continuous parameter has an interior optimum that a one-dimensional
search can locate --- which is exactly what Section~\ref{sec:pending} did, and exactly how it
found that the optimum is the default. A discrete selector has no such notion of
``how much'', and its search space does not admit the same statement.

A matched comparison is left as future work, with the caveat stated above: it would require
budget-matching the two arms and a pre-registered criterion, not a demonstration.

% ------------------------------------------------------------
\section{Discussion}

\subsection{When does arbitration pay?}

The benchmark answers a narrower question than the one it was built to ask, and the narrowing is
the result. \emph{Mixing} pays whenever the rules are complementary: pure Hebbian reaches a Fisher
ratio of 0.472 and pure e-prop 0.331, while their fixed blend reaches 0.651 --- a margin of 0.179
that no single-rule arm approaches. \emph{Arbitration} --- choosing the weight rather than fixing
it --- did not pay here, for a reason that is now specific rather than mysterious: the optimum of
the static mixture is the default value, so the quantity an arbiter would estimate has no room to
move.

This suggests a criterion that can be applied before building an arbiter, in the same syntax as
P2's criterion for when local rules suffice: \textbf{arbitration is worth measuring only if the
static optimum is unknown \emph{and} away from the value one would pick without measuring.} Where
the optimum sits at the default, adaptivity buys nothing and the benchmark cannot distinguish a
good arbiter from a lucky constant.

\subsection{The 1-scalar finding}

Confavreux et al.\ report meta-learning failing as rule parameterisations grow complex, and our
arbiter was deliberately pinned to one scalar to stay on the safe side of that finding
\citep{confavreux2025}. The failure we hit was of a different kind, and worth separating: it was
not that one scalar was too coarse to express the right arbitration --- the static sweep shows the
optimum really is a single number, and that number is 0.5. It was that \emph{no affordable
objective could identify that number}. An integer conflict count saturates; a short-horizon Fisher
probe cannot see the effect it is meant to search; a smooth agreement proxy ranks the endpoints
identically. The bottleneck was not parameterisation size but objective identifiability, which is a
constraint on the outer loop rather than on the arbitration idea.

\subsection{Limits}

All conclusions are limited to a two-rule pair. Three-way mixing was not tested and is out of
scope by design: with more than two rules the definition of conflict itself needs re-design (pairwise
per-row cosines are no longer unique), the number of pairs grows quadratically, and a scheduling
dimension appears. The project roadmap places three-way mixing and scheduling in a later phase.

The benchmark is one task: sequential MNIST under a fixed permutation, 1\% labels, 256 neurons,
five seeds. We did not test other label fractions, other permutations, or larger networks, and the
numbers should not be extrapolated to them. The slow-adaptive arm additionally reserves 10\% of the
training split for validation; its control arm uses the same split, so the comparison is internal,
but it is not directly comparable to the full-data static runs.

Finally, nothing here was run on neuromorphic hardware; the locality claim concerns the learning
rule, not the substrate.
                        % 两规则对 -> 三规则未验; 合成->真实任务

% ------------------------------------------------------------
\section{Conclusion}

Rules conflict on shared synapses, and the conflict is measurable: a fixed mixture of two
heterogeneous local rules separates representations far better than either rule alone
(0.651 against 0.472 and 0.331), in a network with no backpropagation anywhere, under a
protocol that pre-registers its criteria and voids runs whose weights did not move.

Arbitration itself, under the two-sided retention constraint and a one-scalar parameterisation,
came back negative on this benchmark --- measurably negative, not merely untested: the static
optimum is the default value, and the one adaptive scheme that could run without an unaffordable
objective drifted away from that optimum and parked on the constraint boundary. We report the
negative result with the same care as the positive one, because the reusable part is the protocol
that produced it: four degeneracies caught by their own guards, and a criterion that can be
checked before an arbiter is built.

The next step is the cognitive-core integration on the project roadmap, and --- for this line
specifically --- an objective that is both affordable and discriminating.

% ------------------------------------------------------------
\section*{Use of AI assistance}
% 继承 P2 的四段式披露; 更新具体分工 (本论文中 agent 的角色).

% ------------------------------------------------------------
\section*{Data and code availability}
% 仓库 + A/B 记录 DOI + 本论文 Zenodo 存款 DOI (存款后补).

% ------------------------------------------------------------
\bibliographystyle{plainnat}
\bibliography{refs}   % \todo{从 paper/refs.bib 复制所需条目 + Turcu/Confavreux/Zanardi
% 已在台账; 新引用的 Science 2025 多规则突触论文需先核验进台账}

\appendix
\section{Full configuration dumps}      % 继承 P2 附录 A 的格式
\section{Pre-registration records}        % 三个预注册 (参数化扩展/协同基准单方基线/
% 局部性代价接口) 的逐字记录与时间戳
\section{Triviality analysis details}     % 本论文新发现的退化解 (若有) + 继承的两个

\end{document}
